\clearpage\begingroup\SourceLeaf{29}\HistoricalHead{24}{L}{PHYSIK. BRIEFWECHSEL.}
\begin{center}
[2.]\\[.7em]
\textsc{Gauss} an \textsc{Bessel}. Göttingen, 1. April 1827.
\par\medskip\rule{.78\textwidth}{.35pt}\par
{\small [Briefwechsel zwischen \textsc{Gauss} und \textsc{Bessel}, Berlin 1880, S. 470.]}\par
\rule{.78\textwidth}{.35pt}
\end{center}
\noindent .\ .\ .\ .\ .\ .\ .\ .\ .\ .\ .\ .\ .\ .\ .\ .\ .\ .\ .\ .\ .
\par\noindent\hspace*{\parindent}Mit ungemein grossem Interesse habe ich die mir mitgetheilten Nachrichten
über Ihre Pendelversuche gelesen. Wie gross ist doch die Schwierigkeit,
sich der absoluten Grösse einer Constante, die die Naturphaenomene
enthalten, zu versichern! Solange wir dies nicht können, nemlich nicht mit
derselben Schärfe können, scheint es mir höchst unpassend, ein Naturmaass
\textls[120]{ins Leben einführen zu wollen}, und \textls[120]{so lange müssen} alle unsere Maasse
von einem Irgendwo deponirten Individuum abhängig bleiben. Der \textls[120]{französische}
Meter \textls[120]{sollte} ein Naturmaass seyn; allein er hörte auf, es zu seyn,
sobald man festsetzte, er solle \(493^{\mathrm{l}}{,}296\) etc. seyn; von dem Augenblick an
war er ein eben so willkürliches Maass, wie jedes andere, man hat dadurch
die Zahl der willkürlichen Maasse bloss um eins vermehrt, und man hätte
voraussehen können, dass die Zeit kommen würde (wie in \textsc{Puissant} \textit{Instruction
sur l’usage des tables de projection}, 1821, pag. 47), der Erdquadrant sey \textls[120]{nicht}
10000000 Meter sondern 10000724 Meter. Um so grössere Sorgfalt sollte
man aber auf die Conservation der gebrauchten Maasse bei Gradmessungen
und Pendellängen wenden.
\par\noindent .\ .\ .\ .\ .\hspace{.75em} .\ .\ .\ .\ .\ .\ .\ .\ .\ .\ .\ .\ .\ .\ .
\clearpage\endgroup
